\documentclass[10pt,a4paper]{amsart}
\usepackage[margin=19mm]{geometry}
\usepackage{amsmath,amssymb,mathtools}
\usepackage[scr=rsfs,cal=euler]{mathalfa}
\usepackage{xfrac}
\usepackage[unicode,hidelinks,bookmarks=false]{hyperref}
\newcommand{\ie}{i.e.\ }
\newcommand{\cf}{cf.\ }
\newcommand{\resp}{respectively\ }
\newcommand{\Subsection}{\subsection}
\providecommand{\nobreakdash}{\nobreak\hbox{-}}
\setlength{\emergencystretch}{2em}
\multlinegap0pt
\usepackage{mathtools}
\usepackage{amssymb}
\usepackage{enumerate}
\usepackage{tikz}
\usepackage{algorithm}\usepackage{algpseudocode}
\newtheorem{thm}{Theorem}
\newtheorem{lem}{Lemma}
\newcommand{\tr}{\triangleright}
\title{Finite model theory for pseudovarieties and universal algebra: preservation, definability and complexity.}
\author{Lucy Ham, Marcel Jackson}
\date{Source: arXiv:2212.02653v4 (CC BY 4.0)}
\begin{document}
\maketitle

\noindent\emph{Source and licence.} Finite model theory for pseudovarieties and universal algebra: preservation, definability and complexity.. Version arXiv:2212.02653v4 (CC BY 4.0), distributed by arXiv under the Creative Commons Attribution 4.0 International licence, http://creativecommons.org/licenses/by/4.0/, as recorded in the arXiv licence metadata for this exact version and shown on the abstract page https://arxiv.org/abs/2212.02653v4. Full licence notice: \texttt{licenses/arxiv-2212.02653-CC-BY-4.0.txt}.\par\smallskip

\noindent\textit{Editorial scope.} Counted unit: the article's thm thm:FO (source0.tex lines 951--953). The article's complete original proof of it is delivered with it (source0.tex lines 954--964, one \texttt{proof} environment of 11 lines). No mathematics is written, repaired or re-derived: the statement and the proof are the article's own lines, in the article's own notation, and no retained line is substituted or re-flowed.  Retained uncounted as the source-local context the proof needs: lines 543--545; lines 857--859; lines 934--944.  Nothing was omitted from any retained span, so the package carries no omission record.  No retained line contains a citation, so the wrapper carries no bibliography and no bibliography entry.  No retained line contains a figure environment, an image-inclusion command or drawing code, so no image is delivered and the package declares no asset dependency.  Editorial formatting only, in addition: an AMS \texttt{amsart} preamble that restates the article's own declarations and macros used by the retained lines, namely the environments \texttt{thm}, \texttt{lem} on separate counters, together with the standard packages those lines need.  The retained environments print in this document as Theorem 1, Lemma 1 in order of appearance, whereas the article prints the same items with the numbering of its own full document.  The complete native source of the article as distributed in the arXiv e-print for this exact version is bundled unchanged as \texttt{sources/134512/source0.tex}.
\begin{thm}\label{thm:fodecomp}
If $\mathscr{K}$ is a class of algebras of the same signature, closed under taking homomorphic images and with definable cmi congruences, then there is a first order query $\phi(x,y)$ with two free variables $x,y$, such that for every ${\bf A}\in \mathscr{K}$ and every $a\neq b$ in $A$ we have $\phi(a,b)$ defines a subdirectly irreducible quotient of ${\bf A}$ in which the monolith  is generated by representatives of the congruence classes of $a$ and $b$.
\end{thm}

\begin{thm}\label{thm:FOq}
If $\{{\bf L}_1,\dots,{\bf L}_k\}$ is a finite set of semilattice-based finite algebras, then the class $\mathsf{SP}_{\rm fin}^+(\{{\bf L}_1,\dots,{\bf L}_k\})$ is the class of finite models of a $\forall^*\exists^*\forall^*$ sentence in first order logic.  The same is true for $\mathsf{SP}_{\rm fin}(\{{\bf L}_1,\dots,{\bf L}_k\})$.
\end{thm}

\begin{lem}\label{lem:si}
The pointed semidiscriminator variety with $0$ over any finite signature $\mathscr{L}$ has definable cmi congruences.  More precisely, if $\pi_1(x,y,u,v)$ denotes the formula
\[
[((u\wedge v)\tr x\approx x) \vee (u\tr x\not\approx x\And v\tr x\not\approx x)],
\]
then the formula 
\[
\pi_{x,y}(u,v)\coloneqq x\approx y\vee (x\not\approx x\wedge y\And\pi_1(x,y,u,v))\vee(y\not\approx x\wedge y\And\pi_1(x,y,v,u))
\]
defines cmi congruences in the variety of pointed semi-discriminator algebras over~$\mathscr{L}$.
\end{lem}

\begin{thm}\label{thm:FO}
Let ${\bf L}$ be a finite semilattice based algebra with a projection term.  Then the pseudovariety generated by $\flat({\bf L})$ with distinguished $0$ is definable by a $\forall^*\exists^*\forall^*$ sentence of first order logic.
\end{thm}

\begin{proof}
The subdirectly irreducible members of the pseudovariety of $\flat({\bf L})$ are just the flat extensions (with $0$ distinguished) of the members of the universal Horn class of ${\bf L}$.  By Theorem~\ref{thm:FOq}, the class $\mathsf{SP}^+_{\rm fin}({\bf L})$ is first order definable by a $\forall^*\exists^*\forall^*$ sentence $\forall x\forall y\ \Phi(x,y)$ where $\Phi(x,y)$ is an $\exists^*\forall^*$ sentence.  The class of flat extensions of members of this class can then be defined in first order logic by asserting that the structure is flat, with bottom element $0$ (this uses only universal quantifiers), and that $\forall x\forall y\ \Phi(x,y)$ holds for all elements that are not $0$: for every variable $z$ in $\forall x\forall y\ \Phi(x,y)$ we must include $\neg(z\approx 0)$.  (This is where having $0$ is a constant is required in our proof; otherwise it would be necessary to existentially quantify the existence of a zero element, adding an extra layer of quantifier alternation to the sentence.)  Thus we have a sentence $\Psi$ defining the subdirectly irreducible members of
$\mathsf{HSP}_{\rm fin}(\flat({\bf L}))$.
We could now try to apply the one-to-many first order reduction of Theorem~\ref{thm:fodecomp}, but this calls on a linear order, which is acceptable from a complexity-theoretic perspective, but steps outside of pure first order logic of the algebras in consideration.  We can avoid this entirely using the formula $\pi_{x,y}(u,v)$ of Lemma~\ref{lem:si}.  
Let~$\pi_{a,b}$ denote the congruence relation defined by $\pi_{a,b}(u,v)$.  For a model ${\bf K}$ we wish to assert that whenever $a,b\in K$ have $a\neq b$ then ${\bf K}/\pi_{a,b}\models \Psi$.  
Let $\Psi^\flat(a,b)$ denote the result of replacing every instance of equality $u\approx v$ in $\Psi$ by $\pi_{a,b}(u,v)$ (here $u$ and $v$ refer to any terms related by equality in $\Psi$).  Then ${\bf K}/\pi_{a,b}\models \Psi$ for every $a\neq b$ is equivalent to 
\[
{\bf K}\models \forall a\forall b\ \neg(a\approx b)\rightarrow \Psi^\flat(a,b).
\]
Thus the sentence $\forall a\forall b\ \neg(a\approx b)\rightarrow \Psi^\flat(a,b)$ along with the equational axioms defining pointed semidiscriminator varieties with $0$, define $\mathsf{HSP}_{\rm fin}(\flat({\bf L}))$ amongst finite algebras.
\end{proof}

\end{document}
